\scanpage{154}
\begin{center}- 113 -\end{center}
\margin{IX}

9.4.7.3. Cet énoncé garde un sens et reste valable lorsqu'on suppose seulement que $S$ est un schéma noethérien régulier connexe de dimension 1 : au lieu d'un seul accouplement (9.4.7.2), il faut alors considérer un accouplement $Q_s\otimes Q_s^*\longrightarrow\Lambda_s$ pour chaque point fermé $s$ de $S$. C'est sous cette forme plus générale qu'il nous sera le plus commode de prouver 9.4.7.

a) Pour toute extension panachée $X$, l'application
\[
\operatorname{End}(X)\longrightarrow\operatorname{End}(X_\eta)
\]
est bijective. Comme les deux membres sont respectivement les sections de $\uline{\operatorname{Hom}}(P,Q)=D(Q\otimes Q^*)$ sur $S$ et sur $\eta$, l'assertion provient du fait que $S$ est normal et que $D(Q\otimes Q^*)$ est fini sur $S$. De ceci résulte aussitôt que si $X$ et $Y$ sont deux extensions panachées, et si on sait déjà qu'elles sont isomorphes, alors
\[
\uline{\operatorname{Hom}}(X,Y)\longrightarrow\uline{\operatorname{Hom}}(X_\eta,Y_\eta)
\]
est bijectif.

b) Si $X$ et $Y$ sont deux extensions panachées, tout isomorphisme $X_\eta\xrightarrow{\sim}Y_\eta$ provient d'un isomorphisme $X\longrightarrow Y$. Grâce à a) la question est locale pour la topologie étale, ce qui nous permet de nous ramener au cas où $S$ est strictement local. Mais dans ce cas 9.4.3 nous montre que l'application
\margin{\unsure{ON}}
\[
\operatorname{Ext}^1_\Lambda(P,Q)\longrightarrow\operatorname{Ext}^1_\Lambda(P_\eta,Q_\eta)
\]
s'identifie à l'application
\[
V^*\otimes\check H_{\unsure{o|0}}\longrightarrow K^*\otimes\check H_{\unsure{o|0}}
\]
\margin{431}
\margin{431}
\margin{153}
